% 394_Korrekturabstaende_En_ch.tex
% Johann Pascher, 23 Oct 2026
% How far apart the fractal corrections lie

\providecommand{\BM}{\textbf{[B]}}
\providecommand{\KM}{\textbf{[K]}}
\providecommand{\SM}{\textbf{[S]}}
\providecommand{\XM}{\textbf{[X]}}
\providecommand{\QM}{\textbf{[Q]}}

\chapter*{How Far Apart the Fractal Corrections Lie}
\addcontentsline{toc}{chapter}{How Far Apart the Fractal Corrections Lie}

\begin{abstract}
The basic form of FFGFT is ratio-based; a fractal correction appears only in the translation into
SI units. As a principle, the corpus uses a single factor for this, $K_{\text{frak}}=1-100\xi=74/75$,
the per-step factor of the scale recursion (Docs.~133, 300, 384). In the individual calculations,
however, different values and residuals stand side by side (R154). This document writes each of them
as $n$ in $K=1-n\,\xi$, compares it with $n=100$ and sorts by origin. Result: of all candidates, only
one is a genuine test of the SI translation, the fine-structure constant; it requires $n=101.25$,
that is $1.25\,\xi$ away from the rule. All other deviations are algebraic consequences, fitted
quantities, other quantities under the same name, or the distance of the lepton ladder from the Koide
form. The lepton residuals are therefore not a recursion of the measurement but an approximation
error of the ladder, which the Koide form predicts without any measured value. The uniform rule is
not refuted; so far it has been tested exactly once.
\end{abstract}

% ============================================================
\section*{Status markers}
% ============================================================
\begin{center}
\begin{tabular}{@{}lp{11cm}@{}}
\textbf{[B]} & Algebraically proved (rearrangement without additional assumption). \\[2pt]
\textbf{[K]} & Numerically confirmed (check script). \\[2pt]
\textbf{[S]} & Assumption, open, or classification without a calculation of its own. \\[2pt]
\textbf{[X]} & Excluded by the calculation. \\[2pt]
\textbf{[Q]} & Taken from the literature, with source. \\
\end{tabular}
\end{center}
Measured values: $m_e$, $m_\mu$, $\alpha$ and $G$ from CODATA 2022 \QM, $m_\tau=1776.93\pm0.09$\,MeV
from PDG (pdgLive, 2026 edition) \QM. The corpus has so far used $m_\tau=1776.86$\,MeV; where this
makes a difference, both values are given.

% ============================================================
\section{Starting point}
% ============================================================

In the basic form, in natural units with $\alpha=1$, a fractal correction has no place: $\xi
E_0^2=1$ is exact, and in ratios every common factor cancels (A040, Docs.~070, 122). A correction
arises only when a quantity of the basic form is compared with a value measured in SI units. For
this the corpus sets a single factor,
\begin{equation}
  K_{\text{frak}}=1-100\,\xi=\frac{74}{75},
\end{equation}
the per-step factor of the scale recursion $\xi_{n+1}=\xi_n(1-100\,\xi_n)$ (Docs.~133, 300). The
reading that the measurer stands inside this recursion, so that the correction is part of every
measured value, requires the same factor to appear wherever a translation takes place. Whether this
is the case was left open by R154: besides $74/75$, the corpus contains $0.986$, a value required by
$\alpha$, a chain residual for $G$, generation-dependent lepton residuals, a hadronic factor and
$e^{-\xi}$.

To make these values comparable, each is written as $n$ in $K=1-n\,\xi$. The rule is $n=100$; the
distance $n-100$ states by how many $\xi$ a value deviates from it.

% ============================================================
\section{All values on one scale}
% ============================================================

\begin{center}
\begin{tabular}{@{}lrrl@{}}
\toprule
\textbf{Quantity} & \textbf{$K$} & \textbf{$n$} & \textbf{Origin}\\
\midrule
Rule $1-100\xi=74/75$ & $0.986667$ & $100$ & Docs.~133, 300, 384\\
multiplicative form $(1-\xi)^{100}$ & $0.986754$ & $99.34$ & R139\\
required by $\alpha$, $\xi\,m_em_\mu\,\alpha^{-1}$ & $0.986501$ & $101.25$ & R135\\
1 MeV reference & $0.986501$ & $101.25$ & Doc.~386\\
$G$ formula, rounded & $0.986$ & $105$ & Doc.~012\\
$G$ formula with chain residual & $1.0124$ & $-93$ & R141, Doc.~383\\
hadronic factor $144/125$ & $1.152$ & $-1140$ & R123\\
$e^{-\xi}$ & $0.999867$ & $1$ & Doc.~037\\
\bottomrule
\end{tabular}
\end{center}
\KM\ In addition there are the lepton residuals $\varepsilon_e=192\,\xi$, $\varepsilon_\mu=152\,\xi$
and $\varepsilon_\tau=73\,\xi$ (Doc.~352). They measure the distance of the bare ladder
$m_i^\circ=r_i\,\xi^{p_i}\,v\,K_{\text{frak}}$ from the measured value and fall with the generation.

Read as a plain list, this looks like a rule that holds only in principle. The numbers, however,
have different origins, and only some of them can test the rule at all.

% ============================================================
\section{What tests the rule and what does not}
% ============================================================

\subsection{Algebraic consequences}

The residual of $-2.6\,\%$ for $G$ in the one-anchor chain (Doc.~383) is not a finding of its own.
With $m_e$ as anchor, $E_P$ lies $+1.34\,\%$ from the measured value, and because $G=\hbar
c^5/E_P^2$ the deviation of $G$ is exactly twice as large \BM. That it seems to resemble two
correction steps follows from the exponent alone.

The 1 MeV reference from Doc.~386 is the same number as the factor required by $\alpha$. From
$E_0^{\text{bare}}=\sqrt{m_em_\mu}$ it follows that $\bigl(E_0^{\text{bare}}\sqrt{\xi/\alpha}\bigr)^2=
\xi\,m_em_\mu\,\alpha^{-1}$ \BM; the residual of $8\times10^{-5}$ in the reference energy is half of
the residual $1.7\times10^{-4}$ in $K$. It is one finding, not two.

\subsection{Fitted quantities}

The $G$ formula $G=\xi^2/(4m_e)\,C_{\text{conv}}\,K_{\text{frak}}$ matches CODATA only because
$C_{\text{conv}}=7.783\times10^{-3}$ contains the residual of the one-anchor chain; the chain alone
gives $7.58\times10^{-3}$ (R141) \KM. Whether the formula uses $0.986$ or $74/75$ shifts $G$ by
$0.068\,\%$, which is smaller than the fitted part. The $G$ formula can therefore neither confirm nor
refute the rule.

The electron residual of the chain, $-1.32\,\%=-99\,\xi$, lies strikingly close to one correction
step. The chain, however, contains the relation $v/E_P=\xi^4/(5\pi)$ with the factor~10 fitted to the
measured value (Docs.~149, 383); with a different anchor (muon or tau) the residual would be a
different one. The closeness to $100\,\xi$ depends on this choice and carries no weight \SM. The same
holds for the absolute level of the $\varepsilon_i$: it depends on the anchor $v$, which is a
convention of the Standard Model via $G_F$ (R112). Only the differences between the generations are
independent of the anchor.

\subsection{Other quantities under the same name}

$144/125$ is the ratio of two ladder coefficients, $r_\mu/r_\tau=(16/5)/(25/9)=144/125$ \BM, and lies
on the other side of one. It is the bridge factor from Doc.~158 (R140), not an SI correction.
$e^{-\xi}$ in Doc.~037 is a single $\xi$ step with a different role; the document itself notes that
the $G$ formula there uses $1-100\xi$.

% ============================================================
\section{The lepton residuals are the distance of the ladder from the Koide form}
% ============================================================

The lepton ladder $m_i\propto r_i\,\xi^{p_i}$ with $r=(4/3,\,16/5,\,25/9)$ and $p=(3/2,\,1,\,2/3)$
gives the ratios
\begin{equation}
  \frac{m_\mu}{m_e}=\frac{12}{5}\,\xi^{-1/2}=207.846,\qquad
  \frac{m_\tau}{m_\mu}=\frac{125}{144}\,\xi^{-1/3}=16.992 .
\end{equation}
In these ratios $v$, $E_P$ and every common correction factor cancel. They lie $+39.10\,\xi$ and
$+77.56\pm0.38\,\xi$ above the measured values \KM. These are the generation differences of the
$\varepsilon_i$; they cannot contain a common factor such as $K_{\text{frak}}$.

The Koide form with the angle $\theta=2/9$ and the amplitude $d/c=\sqrt2$ (Docs.~353, 367),
\begin{equation}
  \sqrt{m_k}\;\propto\;1+\sqrt2\,\cos\!\Bigl(\tfrac29+\tfrac{2\pi k}{3}\Bigr),
\end{equation}
matches the same ratios to $0.07\,\xi$ and $0.16\pm0.38\,\xi$, for the tau therefore within the
measurement uncertainty \KM. Comparing ladder and Koide form directly with each other, without any
measured value, the ladder lies $39.02\,\xi$ and $77.40\,\xi$ above \KM. Up to $0.2\,\xi$, these are
exactly the lepton residuals.

This settles their origin. The ladder is a rational approximation of the Koide form, and the lepton
residuals are its approximation error. The measurement plays no part in this: the error already lies
between two exact constructions of the basic form. This confirms reading~B from R121 (distance
between cosine and power law, Docs.~367, 369) and, for the ratios, answers the question from R113 of
why the $\varepsilon_i$ fall with the generation.

The ratio of the two residuals is $77.40/39.02=1.984$. It is close to two, but it is not two \KM; the
measurement gives $1.984\pm0.010$ and agrees with the Koide value. An integer step law for the
generations cannot be read off from this \XM.

% ============================================================
\section{What remains}
% ============================================================

After removing the algebraic consequences, the fitted quantities, the unrelated quantities and the
ladder residual, a single value remains that tests the rule independently: the fine-structure
constant. It connects the basic form $\xi E_0^2=1$ with the SI value $\alpha^{-1}=137.036$ via
$E_0^2=m_em_\mu/K$ and requires
\begin{equation}
  K_\alpha=\xi\,m_em_\mu\,\alpha^{-1}=0.986501,\qquad n=101.25 .
\end{equation}
\KM\ The distance from the rule is $1.25\,\xi$ or $168$\,ppm. The measurement uncertainty of
$K_\alpha$ is below $3\times10^{-8}$; the residual is therefore real and not explained by measurement
errors. The additive form $1-100\xi$ lies closer than the multiplicative form $(1-\xi)^{100}$
($1.25\,\xi$ against $1.90\,\xi$), but does not separate them more sharply than their mutual distance
(R139). A second-order term, $1-\varepsilon-\varepsilon^2$ with $\varepsilon=1/75$, would give
$n=101.33$ and misses $K_\alpha$ by $12$\,ppm; it is not adopted \SM.

The uniform rule is therefore not refuted. So far it has been tested exactly once and passes this
test at $1.25\,\xi$. Any further statement about its uniformity needs a second quantity that, like
$\alpha$, connects a quantity of the basic form with an SI measured value without containing a
fitted factor. The natural candidate is the muon lifetime $\tau_\mu$ directly from FFGFT (R112) \SM.

% ============================================================
\section{Conclusion}
% ============================================================

The list from R154 looked like many different corrections. Sorted by origin, what remains is one
rule with one test. The lepton residuals do not belong to the SI translation but to the structure of
the basic form itself: they disappear when the ladder is replaced by the Koide form. The reading
that the measurer stands inside the recursion therefore concerns only the translation into SI values,
and there so far only $\alpha$. The ratios need no correction; in the Koide form they are exact up to
the measurement accuracy.

What remains open is the residual of $1.25\,\xi$ for $\alpha$ and the form of the factor (R135,
R139), a second independent test of the rule (R112), and the origin of the factor~10 in $v/E_P$
(Docs.~149, 387) \SM.

\medskip
\noindent\textbf{Check script:} \path{2/python/Dok394_Skripte/pruef_394_korrekturabstaende.py}
